%%%%%%%%%%%%%%%%%%%%%%%%%%%%%%%%%%%%%%%%%%%%%%%%%%%%%%%%%%%%%%%%%%%%%%
%%                       IAEW LaTeX Template                        %%
%%%%%%%%%%%%%%%%%%%%%%%%%%%%%%%%%%%%%%%%%%%%%%%%%%%%%%%%%%%%%%%%%%%%%%

\chapter{Zusammenfassung und Ausblick}
\label{ch:conc}
% ===============================================================
% GESCHRIEBEN am 26.09.2026 nach dem Absatzplan, den der Verfasser am
% selben Tag freigegeben hat: zwoelf Absaetze zu je acht Saetzen, drei
% Seiten. Vorgehen: alle Aussagen der Kapitel 2 bis 5 gezogen, auf fuenf
% Kernpunkte verdichtet, je Kernpunkt ein bis zwei Absaetze. Die fuenf
% Kernpunkte sind so gefasst, dass sie ohne die Zahlen als Behauptung
% stehen; die Zahlen selbst stehen in Kapitel 4, hier nur verbale
% Groessenordnungen wie in Kapitel 5.
%   K1  Die kurative Vorhaltung laesst sich als Marktprodukt ausgestalten.
%   K2  Der Zuschnitt folgt aus den Anforderungen, jede Anforderung hat
%       ihren Preis.
%   K3  Der kurative Reservierungspreis ist ein Preis gegen die
%       Regelleistung, und er ist bezahlbar.
%   K4  BESS muessen in das Engpassmanagement, kurativ passt zu ihrer
%       Technologie: praeventiv verlangt Energie ueber die Dauer des
%       Engpasses, kurativ schnelle Leistung fuer ein kurzes Intervall.
%       Vorgabe des Verfassers; nicht gesetzt ist, der praeventive
%       Redispatch lasse sich mit BESS nicht durchfuehren, denn
%       chapter_2.tex Zeile 1068 sagt, der Energieinhalt begrenze die
%       Eignung.
%   K5  Kurativ braucht einen eigenen Rahmen; systemweit wird aus der
%       Sicherheits- eine Effizienzmassnahme.
% Keine Zitate, alles ist in den Kapiteln 2 bis 5 belegt. Keine
% Doppelpunkte, Gedankenstriche, Semikola. Ein Hauptsatz und hoechstens
% ein Nebensatz je Satz. Kein Pronomen ueber die Satzgrenze.
% Der Kommentarkopf der Vorfassung vom 18.09.2026 steht darunter, die
% alte Punkteliste fuer den Ausblick am Ende der Datei als Kommentar.
% ===============================================================
% VORFASSUNG, Kommentarkopf vom 18.09.2026, unveraendert:
% Zielumfang Kapitel 6: 5 Seiten
%   6.1 Zusammenfassung entlang der Forschungsfragen ......... 3
%   6.2 Ausblick ............................................. 2
% ACHTUNG 18.09.2026: Der Verfasser hat die drei Teilfragen in Abschnitt
% 1.2 gestrichen, sie stehen dort nicht mehr und sollen nirgends wieder
% auftauchen.
% ZWEI BEFUNDE, die in der Zusammenfassung stehen muessen: die
% Verguetungsstruktur der geltenden Festlegung ist nicht additiv; die
% Verlagerung des Anreizproblems von der Menge des Eingriffs auf den
% Nachweis der Verfuegbarkeit. Beide stehen jetzt in 6.1 Absatz 3 und
% 6.2 Absatz 1.
% ===============================================================

Das Engpassmanagement im deutschen Übertragungsnetz ist präventiv angelegt und kostet jedes Jahr Milliarden, denn der \ac{ÜNB} hält den Grenzwert jedes Betriebsmittels schon vor dem Fehlerfall ein.
Die kurative Systemführung nutzt dagegen die Überlast, die ein Betriebsmittel nach dem Fehler für kurze Zeit trägt, und lässt das Netz davor höher auslasten.
Dafür braucht der \ac{ÜNB} Anlagen, die nach dem Fehler innerhalb der zulässigen Zeit ihre Leistung ändern, und diese Bereitschaft heißt kurative Vorhaltung.
Eine Vorhaltung kann der \ac{ÜNB} nicht anordnen, und das Engpassmanagement vergütet sie bislang nicht.
Die vorliegende Arbeit entwirft deshalb die kurative Reservierung als Produkt, mit dem der \ac{ÜNB} die Vorhaltung marktlich beschafft, und bestimmt den kurativen Reservierungspreis, den ein Betreiber dafür mindestens fordern muss.
Gerechnet ist der Preis für ein \ac{BESS} aus der Sicht seines Betreibers, mit den beobachteten Preisen des Jahres 2025 an den kurzfristigen Strommärkten und in der Regelleistung.
Eingeordnet wird der Entwurf für das Trendjahr des \ac{NEP}, in dem Wind, Photovoltaik und Speicher den Hauptanteil der Leistung stellen.
Fünf Erkenntnisse tragen die Arbeit, und im Folgenden werden sie benannt und auf die kurative Systemführung bezogen.

% ---------------------------------------------------------------
\section{Zusammenfassung}
\label{sec:summary}

Die erste Erkenntnis lautet, dass sich die kurative Vorhaltung als Marktprodukt ausgestalten lässt.
Die kurative Reservierung bindet je Stunde ein Leistungsband und bei Speichern ein Ladezustandsband, das der Betreiber über die Bindungsdauer freihält.
Vergütet wird die Vorhaltung und nicht der Eingriff, denn die Kosten des Betreibers entstehen aus dem Freihalten der Leistung und nicht erst aus einem Abruf.
Ein Abruf wird gesondert vergütet, nach der gelieferten Energie und der Nachlaufpflicht, sodass der kurative Reservierungspreis frei von der Abrufwahrscheinlichkeit bleibt.
Die Zusage bindet über eine Pönale und über den Nachweis, dass die angebotene Leistung zum Zeitpunkt der Einplanung frei ist.
An den acht Anforderungen, die aus der kurativen Systemführung folgen, hält der Entwurf sieben am Produkt ein, nämlich Reaktionszeit, Verbindlichkeit, Lokationalität, Bemessungsgegenstand, Diskriminierungsfreiheit, Teilnahme und Verträglichkeit.
Die achte, die Umsetzbarkeit, entscheidet sich nicht am Produkt, sondern am Rahmen und am Betrieb des \ac{ÜNB}.
Ein Marktprodukt für die kurative Vorhaltung ist damit kein offenes Konzept mehr, sondern ein Entwurf, dessen Lücken sich benennen lassen.

Die zweite Erkenntnis lautet, dass der Zuschnitt des Produkts aus den Anforderungen folgt und jede Anforderung ihren Preis hat.
Die drei Reaktionszeitklassen folgen aus der Reaktionszeit, die das betroffene Betriebsmittel vorgibt, und die Vergütung an Leistung und Bindungsdauer folgt aus dem Bemessungsgegenstand.
Die Ausschreibung je Stunde und ihre Lage hinter der Regelleistung folgen aus Teilnahme und Verträglichkeit, denn ein Betreiber vermarktet zuerst die Regelleistung und bietet dann, was offen bleibt.
Ein Preis je Tag statt je Stunde bindet die Stunden unvollständig und kostet dennoch beinahe den vollen Preis, weil er gerade die teuren Stunden frei lässt.
Die vorgehaltene Energie je Abruf bewegt den Preis nur wenig, und über sie entscheidet allein der \ac{ÜNB} nach der Zeit bis zur ablösenden Maßnahme.
Eine schlanke und knotenscharfe Präqualifikation folgt aus Diskriminierungsfreiheit und Lokationalität, denn nur ein offener Zugang füllt einen Netzknoten mit genügend Geboten.
Je schärfer eine Anforderung das Produkt schneidet, desto kleiner wird der Kreis der Gebote am Netzknoten und desto höher der Knappheitsaufschlag, den ein Gebot über den ermittelten Preis legt.
Die Anforderungen sind deshalb so hoch wie nötig und so niedrig wie möglich zu setzen, und den Maßstab gibt allein die Sicherheit des Netzbetriebs.

Die dritte Erkenntnis lautet, dass der kurative Reservierungspreis ein Preis gegen die Regelleistung ist und dass er bezahlbar bleibt.
Der Preis misst die Vermarktung, die ein Betreiber für die Reservierung einer Stunde aufgibt, und er entsteht für jede Stunde des Jahres neu.
Sein Niveau trägt der Leistungspreis der \ac{aFRR}, denn die Vorhaltung von \ac{aFRR}-Leistung bindet den größten Teil der Leistung und entfällt mit der reservierten Stunde vollständig.
Die Verteilung ist schief, denn die typische Stunde ist günstig, während wenige Stunden mit hohen Handelsspannen am \ac{IDC} den Durchschnitt tragen.
Über den Tag ist die Ladereservierung mittags teuer und die Entladereservierung am Vormittag und am Abend, sodass eine bilanziell ausgeglichene Vorhaltung nachts am günstigsten ist.
Über das Jahr ist die Reservierung im Winter um rund zwei Fünftel günstiger als im Sommer, und die Ladereservierung bildet von März bis Oktober ein Band über die Mittagsstunden.
Unter vollständiger Bindung zahlt der \ac{ÜNB} ein knappes Viertel mehr, als die Anlage an den Märkten verdient hätte, und unter niedrigeren Marktpreisen sinkt der Preis mit dem Markt.
Der kurative Reservierungspreis hängt damit an einem einzelnen Markt und erbt dessen Entwicklung.

Gegen den präventiven Redispatch gehalten, bleibt der kurative Reservierungspreis im Tagesmittel fast das ganze Jahr unter dessen Arbeitspreis.
Der Vergleich trägt eine Unsicherheit, denn der Arbeitspreis vergütet bewegte Arbeit und ist je Technologie kostenbasiert bemessen, der kurative Reservierungspreis dagegen die Vorhaltung.
Kurative Maßnahmen lassen sich damit in einem großen Teil der Stunden wirtschaftlich einplanen, nicht aber in jeder.
Der größte Bedarf an Engpassmanagement fällt in den Winter, wo die Reservierung am günstigsten ist, sodass Bedarf und Preis zusammenfallen.
Die Ausnahme ist der Mittag, denn dort trifft der größte Bedarf an Reduzierung der Einspeisung auf die teuerste Ladereservierung.
Eine Maßnahme aus abgeregelter Photovoltaik und entladendem \ac{BESS} bleibt dort dennoch günstig, weil die Entladereservierung zur selben Zeit wenig kostet.
Wind- und Photovoltaikanlagen bieten die Rücknahme der Einspeisung voraussichtlich günstiger an, wirken aber nur in eine Richtung und nur bei passendem Wetter.
Ob sich die kurative Vorhaltung lohnt, entscheidet sich deshalb nicht am Preis allein, sondern im Verhältnis von Preis und Aufwand der Systemführung, und beides bestimmt der \ac{ÜNB} mit.

Die vierte Erkenntnis lautet, dass \ac{BESS} in das Engpassmanagement eintreten müssen und dass die kurative Systemführung der Weg ist, der zu ihrer Technologie passt.
Ein \ac{BESS} bewegt große Leistung kurzfristig und handelt allein nach Marktsignalen, sodass es einen Engpass verschärfen kann, ohne es zu wissen.
Der Betriebsplanung des Engpassmanagements entziehen sich damit viele Leistungsänderungen, und der wachsende Bestand stößt am Netzanschluss an Grenzen.
Der \ac{ÜNB} setzt ein \ac{BESS} im Engpassmanagement bislang nicht ein, denn die kostenbasierte Vergütung des Redispatch trifft den Werteverbrauch und die Opportunität eines Speichers nicht.
Der präventive Redispatch verlangt vom Speicher zudem, was der Speicher nicht hat, nämlich Energie über die ganze Dauer des Engpasses.
Die kurative Systemführung verlangt, was ein Speicher hat, nämlich schnelle Leistung in beide Richtungen für ein kurzes Intervall bis zur ablösenden Maßnahme.
Die kurative Systemführung wirkt damit in zwei Hinsichten, denn sie bindet den Speicher in das Engpassmanagement ein und macht seine Technologie netzdienlich nutzbar.
Ein Zuschlag sagt dem Betreiber zudem, wo seine Anlage dem Netz nützt, und eine solche Auskunft gibt ihm kein anderer Markt.

Die fünfte Erkenntnis lautet, dass die kurative Systemführung einen eigenen Rahmen braucht und systemweit aus der Sicherheits- eine Effizienzmaßnahme wird.
Eine Vergütung der vorgehaltenen Leistung kann neben der Vergütung der bewegten Arbeit bestehen, sofern die Arbeit den Engpass dauerhaft behebt und die Leistung das Überlastintervall überbrückt.
Eine solche Vergütung lässt sich nicht in das Duldungsregime des Redispatch einfügen, sodass §\,13a \ac{EnWG} zu erweitern ist.
Der Redispatch ändert einen Fahrplan einmalig, während die kurative Reservierung ihn vorab über die Bindungsdauer einschränkt.
Solange der Markt die Übertragungskapazität nicht selbst abbildet, übernimmt der kurative Marktmechanismus diese Aufgabe bei geringem Eingriff in die Marktergebnisse.
Mit einer Ausrollung an vielen Netzknoten nutzt die kurative Maßnahme die bestehende Kapazität planmäßig aus, sofern ein Bestand wirksamer Anlagen je Netzknoten und ein gleichmäßig belastbares Netz bestehen.
Begrenzungen am Netzanschluss, Vorgaben für den Betrieb und ein Gebotszonensplit greifen am Verhalten des Akteurs an und sind damit Übergangslösungen.
Mit vollständiger Einbindung wird mehr Speicherleistung am bestehenden Netz anschließbar, sodass die kurative Systemführung zu einer Voraussetzung für den Ausbau der Speicher wird.

Diese fünf Erkenntnisse stehen unter den Grenzen der Untersuchung, die ihre Lesart bestimmen.
Das Optimierungsmodell rechnet eine einzelne Anlage und bildet weder das Netz noch den Wettbewerb mehrerer Anbieter noch die Rückwirkung der Speicher auf die Preise ab.
Das Optimierungsmodell kennt alle Preise des Liefertages, sodass der Optionswert der Bindung entfällt und der Preis die Opportunität nicht unterschätzt.
Die Preissuche findet ein koordinatenweises Minimum, sodass der ausgewiesene Preis der vorsichtigere Wert ist.
Die Wahrscheinlichkeit eines Abrufs und die Pönale bleiben außer Betracht, und der Leistungspreis der \ac{aFRR} ist wegen seines Gebotsschlusses Rechengröße und nicht Opportunität.
Der ermittelte kurative Reservierungspreis ist deshalb der opportunitätskostenbasierte Teil eines Gebots und nicht das Gebot.
Ein Gebot liegt darüber, um den Aufwand des Anbietens, das Risiko aus Abruf und Pönale und die Knappheit am Netzknoten.
Wer die Erkenntnisse nutzt, liest sie als Aussagen über Niveau, Muster und Einflüsse des Preises und nicht als Prognose einzelner Stunden.

% ---------------------------------------------------------------
\section{Ausblick}
\label{sec:outlook}

Damit sich die kurative Systemführung durchsetzt, muss zuerst der Rahmen entstehen, den es heute nicht gibt.
§\,13a \ac{EnWG} ist um eine Vergütung der vorgehaltenen Leistung zu erweitern, denn das Duldungsregime trägt eine Vorhaltung nicht.
Die Vergütungsstruktur muss Vorhaltung und Abruf trennen und für den Abruf das Indifferenzprinzip des Redispatch übernehmen.
Ein Verfahren zum Nachweis der Verfügbarkeit gehört in den Kern des kurativen Marktmechanismus, denn der Anreiz verschiebt sich von der Menge des Eingriffs auf die Zusage freier Leistung.
Die Präqualifikation kann auf der Regelleistung aufsetzen, sodass allein der Nachweis der kurativen Verfügbarkeit hinzukommt.
Zu vermeiden ist ein strategisches Bietverhalten, denn ein Knappheitsaufschlag treibt die Kosten der Vorhaltung, ohne eine Leistung zu vergüten.
Dagegen hilft ein offener Zugang, der die Zahl der Anbieter je Netzknoten hebt.
Ein Gebotszonensplit veränderte das Engpassmanagement, sodass der kurative Reservierungspreis neu zu bewerten wäre.

Neben dem Rahmen entscheidet der Betrieb, denn eine kurative Maßnahme trägt nur mit ihrer Wirkungskette aus Fehlererkennung, Signalübertragung und Anlagenreaktion.
Der Ausfall eines Gliedes der Wirkungskette wirkt wie der Ausfall eines Betriebsmittels, sodass eine Redundanz in Fehlerdetektion, Anlagensteuerung und Maßnahmenabruf Voraussetzung ist.
Die Ablösung der kurativen Maßnahme ist festzulegen, denn die Anlaufzeit der Ablösung bestimmt die vorgehaltene Energie und damit den Preis.
Ein zusätzliches Produkt mit längerem Zeitfenster übernähme die Ablösung im Markt selbst.
Das Anschlussverfahren ist auf eine Technologie mit hoher Leistung und geringem Energieinhalt zuzuschneiden, statt sie am Netzanschluss zu begrenzen.
Ein Bestand wirksamer Anlagen je Netzknoten und ein Prozess, der die Maßnahmen verlässlich einplant, sind heute nicht gegeben.
Der Übergang von der Sicherheits- zur Effizienzmaßnahme ist deshalb ein Entwicklungspfad, den der \ac{ÜNB} mit dem Netzausbau und mit den Anschlussregeln gestaltet.
Aus Projekten im Echtbetrieb ist zu lernen, wie häufig und wie lange eine kurative Maßnahme abgerufen wird und wie sich die Verfügbarkeit im Betrieb nachweisen lässt.

Am Modell bleibt Forschungsbedarf, der die Erkenntnisse vom Einzelfall auf das System hebt.
Eine systemweite Aussage verlangt ein Modell, das die Abbildung der Märkte mit der Netzsicherheitsrechnung und der Planung des Engpassmanagements verbindet.
Voran steht eine wirtschaftliche Bewertung an einem Datensatz der Netzentwicklungsplanung, die den Einfluss auf das Engpassmanagement quantitativ erfasst.
Die Wahrscheinlichkeit eines Abrufs und damit der Optionswert der Bindung sind zu bestimmen, wofür das Optimierungsmodell um den Abruf als stochastische Größe zu erweitern ist.
Die Menge der Vorhaltung ist auf Seiten des \ac{ÜNB} gegen ein Spektrum möglicher Netzzustände zu bemessen und nicht gegen einen geplanten Betriebspunkt.
Das Bietverhalten und der Wettbewerb mehrerer Anbieter an einem Netzknoten sind abzubilden, ebenso die Vermarktung im Verbund mehrerer Anlagen.
Der Handel am \ac{IDC} ist mit einer rollierenden Bewertung abzubilden, denn der Faktor auf die Handelsspanne trifft den mehrfachen Handel derselben Lieferstunde nicht.
Zu bestimmen ist zudem die zusätzlich anschließbare Speicherleistung, denn darin liegt der Nutzen für das Stromsystem.

Die kurative Systemführung ist mehr als ein weiteres Werkzeug des Netzbetriebs.
Die kurative Systemführung macht die Technologie der Speicher netzdienlich nutzbar und öffnet ihnen das Engpassmanagement, in dem sie bislang nicht vorkommen.
Der Preis dafür ist bezahlbar, im Winter am niedrigsten und über die Jahre an die Regelleistung gekoppelt.
Was fehlt, ist kein Ergebnis mehr, sondern ein Rahmen, ein Betrieb und ein Netz, die die Einbindung zulassen.
Mit Rahmen, Betrieb und Netz wird aus der Sicherheitsmaßnahme eine Maßnahme der Effizienz, und mit der Effizienzmaßnahme wird mehr Speicherleistung am bestehenden Netz anschließbar.
Im Trendjahr des \ac{NEP} stellen Wind, Photovoltaik und Speicher den Hauptanteil der Leistung, und ein Engpassmanagement für dieses System nutzt diese Leistung, statt sie zu begrenzen.
Die kurative Reservierung ist dafür das Produkt, und der kurative Reservierungspreis ist der Maßstab ihrer Beschaffung.
Über die kurative Systemführung entscheidet damit nicht der Preis, sondern der Rahmen, der sie zulässt.

% ===============================================================
% ALTE PUNKTELISTE FUER DEN AUSBLICK, Stand 18.08. bis 20.09.2026,
% am 26.09.2026 durch die Absaetze 9 bis 11 ersetzt. Aufgenommen sind
% Redundanzkonzepte (Absatz 10), stochastische Erweiterung (11),
% Portfolio (11), Marktdesign mit Bietwettbewerb und Beschaffungsmenge
% (11), Bemessung gegen ein Spektrum von Netzzustaenden (11) sowie aus
% AUSBLICK.md die Wirksamkeit je Netzknoten (11), die
% Rolling-Intrinsic-Bewertung (11) und die IDC-Handelsstrategie (11).
% NICHT aufgenommen, Entscheidung des Verfassers vom 26.09.2026: die
% laengere Auswertung nach der Umstellung des Day-Ahead-Marktes, weil
% die Arbeit das ganze Jahr 2025 rechnet, und der Vergleich mit MACSE
% und dem Local Constraint Market, weil die Primaerquellen ungeprueft
% sind. Zurueckgenommene Fassung, nicht wieder aufzunehmen:
%\begin{itemize}
%    \item Redundanzkonzepte für kurative Maßnahmen als Voraussetzung der Betriebsimplementierung, angekündigt in Abschnitt \ref{sec:challenges} \cite{sous_untersuchung_2022, sous_comparison_2025}
%    \item Stochastische Erweiterung des Modells mit der Abrufwahrscheinlichkeit als Szenarioparameter
%    \item Übertragung von der Einzelanlage auf ein Anlagenportfolio
%    \item Ausgestaltung eines umfassenden Marktdesigns für die kurative Systemführung einschließlich Bietwettbewerb und Beschaffungsmenge
%    \item Bemessung der vorzuhaltenden Menge gegen ein Spektrum möglicher Netzzustände anstelle eines geplanten Betriebspunkts
%    \item Auswertung über einen längeren Zeitraum, sobald hinreichend viele Monate nach der Umstellung des Day-Ahead-Marktes vorliegen
%    \item Vergleich mit internationalen Beschaffungsverfahren für lokale Flexibilität als Maßstab für die Ausgestaltung
%\end{itemize}
